% =====================================================================
\chapter{Soluciones de los ejercicios}
\label{ap:soluciones}
% =====================================================================

Cada solución se presenta en cinco partes: datos, planteamiento, desarrollo, resultado y comprobación. En los pasos intermedios se redondea a cuatro cifras significativas salvo que se indique otra cosa; los cálculos con números duales y con enteros son exactos. Todos los resultados numéricos se han comprobado con Python y PyTorch en doble precisión.

% ---------------------------------------------------------------------
\section{Tema 1}

\begin{solucion}{eje:1-traza}
\datos la traza del \cref{ej:baydin}: $v_{-1},v_0$ (entradas) y $v_1,\dots,v_5$.
\planteamiento contar un nodo por variable y una arista por cada par $j\pre i$ (\cref{def:grafo}); un orden topológico solo exige que cada nodo vaya detrás de sus argumentos (\cref{def:topologico}).
\pasos
\begin{enumerate}[label=\arabic*.]
  \item Nodos: $v_{-1},v_0,v_1,\dots,v_5$, en total $7$.
  \item Aristas: $v_1$ tiene un argumento, $v_2$ dos, $v_3$ uno, $v_4$ dos y $v_5$ dos: $1+2+1+2+2=8$.
  \item Sucesores: $v_{-1}$ alimenta a $v_1$ y $v_2$; $v_0$, a $v_2$ y $v_3$. Los demás nodos tienen un solo sucesor, salvo $v_5$, que no tiene ninguno.
  \item Otro orden topológico: $v_0,v_{-1},v_3,v_2,v_1,v_4,v_5$. Cada nodo aparece tras sus argumentos: $v_3$ tras $v_0$; $v_2$ tras $v_{-1}$ y $v_0$; $v_1$ tras $v_{-1}$; $v_4$ tras $v_1$ y $v_2$; $v_5$ tras $v_4$ y $v_3$.
\end{enumerate}
\resultado \textbf{7 nodos y 8 aristas; $v_{-1}$ y $v_0$ tienen dos sucesores}; un orden alternativo es $v_0,v_{-1},v_3,v_2,v_1,v_4,v_5$.
\comprobacion la suma de los grados de entrada de los nodos interiores coincide con el número de aristas: $8$.
\end{solucion}

\begin{solucion}{eje:1-precision}
\datos diferencias centradas con $h=10^{-5}$ y con $h=10^{-15}$, AD; doble precisión, $\uep\approx 1{,}1\cdot 10^{-16}$.
\planteamiento el error de la diferencia centrada es truncamiento, de orden $h^{2}$, más redondeo, de orden $\uep/h$ (módulo~4, tema~6); la AD solo arrastra el redondeo de cada operación.
\pasos
\begin{enumerate}[label=\arabic*.]
  \item $h=10^{-15}$: truncamiento del orden de $10^{-30}$, despreciable; redondeo del orden de $10^{-16}/10^{-15}=10^{-1}$. La resta $f(x+h)-f(x-h)$ pierde casi todas las cifras.
  \item $h=10^{-5}$: truncamiento del orden de $(10^{-5})^{2}=10^{-10}$ y redondeo del orden de $10^{-16}/10^{-5}=10^{-11}$. Error total del orden de $10^{-10}$.
  \item AD: aplica reglas exactas, como $(e^{x})'=e^{x}$, a los valores; el error relativo es del orden de $\uep$.
  \item Con $f=e^{x}$ en $x=1$: el error medido es $5{,}9\cdot 10^{-11}$ con $h=10^{-5}$ y $0{,}17$ con $h=10^{-15}$.
\end{enumerate}
\resultado \textbf{de menor a mayor precisión: $h=10^{-15}$, $h=10^{-5}$, AD.}
\comprobacion los errores medidos siguen el orden previsto: $0{,}17\gg 5{,}9\cdot 10^{-11}\gg 10^{-16}$.
\end{solucion}

\begin{solucion}{eje:1-d2l}
\datos $f(x)=\ln(x^{2})\cdot\sen x+x^{-1}$, $x=2$.
\planteamiento descomponer en operaciones de la \cref{def:elementales}: un nodo por cada potencia, logaritmo, seno, producto y suma.
\pasos
\begin{enumerate}[label=\arabic*.]
  \item (a) Nodos: $a=x^{2}$, $b=\ln a$, $c=\sen x$, $d=b\cdot c$, $e=x^{-1}$, $f=d+e$. Aristas: $x\to a\to b\to d$, $x\to c\to d$, $x\to e$, $d\to f$, $e\to f$. La entrada $x$ tiene tres sucesores: $a$, $c$ y $e$.
  \item (b) $a=2^{2}=4$.
  \item $b=\ln 4=1{,}386$.
  \item $c=\sen 2=0{,}9093$.
  \item $d=1{,}386\cdot 0{,}9093=1{,}261$.
  \item $e=1/2=0{,}5$.
  \item $f=1{,}261+0{,}5=1{,}761$.
\end{enumerate}
\resultado \textbf{grafo con seis nodos interiores y $f(2)\approx 1{,}761$.}
\comprobacion evaluando directamente, $\ln 4\cdot\sen 2+0{,}5=1{,}2606+0{,}5=1{,}7606$.
\end{solucion}

\begin{solucion}{eje:1-complejo}
\datos $f(x)=e^{x}$; $h=10^{-1}$ y $h=10^{-8}$; $x=1$ para las cifras.
\planteamiento usar $e^{x+ih}=e^{x}(\cos h+i\sen h)$ y comparar con la diferencia centrada $\bigl(e^{x+h}-e^{x-h}\bigr)/(2h)=e^{x}\sinh(h)/h$.
\pasos
\begin{enumerate}[label=\arabic*.]
  \item $\operatorname{Im}e^{x+ih}=e^{x}\sen h$, luego $\operatorname{Im}f(x+ih)/h=e^{x}\,\sen(h)/h$.
  \item Error relativo respecto a $f'(x)=e^{x}$: $1-\sen(h)/h$. Con $\sen h=h-h^{3}/6+O(h^{5})$, es $h^{2}/6+O(h^{4})$.
  \item $h=10^{-1}$: paso complejo $1-\sen(0{,}1)/0{,}1=1{,}666\cdot 10^{-3}$; centrada $\sinh(0{,}1)/0{,}1-1=1{,}668\cdot 10^{-3}$. Domina el truncamiento y ambos son casi iguales.
  \item $h=10^{-8}$: paso complejo, $h^{2}/6\approx 1{,}7\cdot 10^{-17}$, por debajo de $\uep$; el valor calculado coincide con $e$ en todas las cifras. Centrada: truncamiento $1{,}7\cdot 10^{-17}$, pero el redondeo, de orden $\uep/h\approx 10^{-8}$, da un error medido de $2{,}4\cdot 10^{-9}$.
\end{enumerate}
\resultado \textbf{error relativo $1-\sen(h)/h\approx h^{2}/6$}: $1{,}67\cdot 10^{-3}$ con $h=0{,}1$ en ambos métodos; con $h=10^{-8}$, error nulo en doble precisión con paso complejo y $2{,}4\cdot 10^{-9}$ con la centrada.
\comprobacion el paso complejo no resta cantidades próximas, de modo que su error decrece con $h$ sin límite (\cref{obs:paso-complejo}); la centrada alcanza su mínimo hacia $h\approx 10^{-5}$.
\end{solucion}

\begin{solucion}{eje:1-logistica}
\datos $l_1=x=0{,}2$, $l_1'=1$; $l_{k+1}=4l_k(1-l_k)$, $l_{k+1}'=4l_k'(1-2l_k)$.
\planteamiento propagar el par (valor, derivada) paso a paso, como en el modo directo.
\pasos
\begin{enumerate}[label=\arabic*.]
  \item $l_2=4\cdot 0{,}2\cdot 0{,}8=0{,}64$; $l_2'=4\cdot 1\cdot(1-0{,}4)=2{,}4$.
  \item $l_3=4\cdot 0{,}64\cdot 0{,}36=0{,}9216$; $l_3'=4\cdot 2{,}4\cdot(1-1{,}28)=9{,}6\cdot(-0{,}28)=-2{,}688$.
  \item Comprobación simbólica: $l_3=16(x-x^{2})(1-2x)^{2}$, y por la regla del producto
  $l_3'=16\bigl[(1-2x)^{3}-4(x-x^{2})(1-2x)\bigr]=16(1-2x)(1-8x+8x^{2})$.
  \item En $x=0{,}2$: $16\cdot 0{,}6\cdot(1-1{,}6+0{,}32)=9{,}6\cdot(-0{,}28)=-2{,}688$.
\end{enumerate}
\resultado \textbf{$l_2=0{,}64$, $l_2'=2{,}4$, $l_3=0{,}9216$, $l_3'=-2{,}688$.}
\comprobacion $l_3=16\cdot 0{,}2\cdot 0{,}8\cdot 0{,}36=0{,}9216$ coincide con el valor de la recurrencia.
\end{solucion}

\begin{solucion}{eje:1-grado}
\datos $l_1=x$, $l_{k+1}=4l_k-4l_k^{2}$.
\planteamiento inducción sobre $k$ para el grado; recuento directo de operaciones para el par $(l_k,l_k')$.
\pasos
\begin{enumerate}[label=\arabic*.]
  \item Base: $l_1=x$ tiene grado $1=2^{0}$ y coeficiente principal $c_1=1\ne 0$.
  \item Hipótesis: $l_k$ tiene grado $2^{k-1}$ y coeficiente principal $c_k\ne 0$.
  \item Entonces $l_k^{2}$ tiene grado $2^{k}$ y coeficiente principal $c_k^{2}$, mientras que $4l_k$ tiene grado $2^{k-1}<2^{k}$. El término de grado $2^{k}$ de $l_{k+1}$ es $-4c_k^{2}x^{2^{k}}$, con $c_{k+1}=-4c_k^{2}\ne 0$: grado exactamente $2^{k}$.
  \item Un polinomio de grado $2^{k-1}$ tiene $2^{k-1}+1$ coeficientes, de modo que la forma desarrollada de $l_k$, y con ella la de $l_k'$, crece exponencialmente con $k$.
  \item Par $(l_{k+1},l_{k+1}')$ a partir de $(l_k,l_k')$: $t=1-l_k$ (1 operación); $l_{k+1}=4\cdot l_k\cdot t$ (2); $s=t-l_k=1-2l_k$ (1); $l_{k+1}'=4\cdot l_k'\cdot s$ (2). Total: 6 operaciones por paso.
  \item De $(l_1,l_1')=(x,1)$ a $(l_k,l_k')$ hay $k-1$ pasos: a lo sumo $6(k-1)$ operaciones.
\end{enumerate}
\resultado \textbf{$\deg l_k=2^{k-1}$, mientras que el modo directo calcula $(l_k,l_k')$ con $6(k-1)$ operaciones}: coste lineal frente a una expresión de tamaño exponencial.
\comprobacion $l_3=16x(1-x)(1-2x)^{2}$ tiene grado $1+1+2=4=2^{2}$, con coeficiente principal $16\cdot(-1)\cdot 4=-64=-4c_2^{2}$, pues $c_2=-4$.
\end{solucion}

% ---------------------------------------------------------------------
\section{Tema 2}

\begin{solucion}{eje:2-cuadrado}
\datos $(3+\varepsilon)^{2}$; $f\colon\R^{1000}\to\R$.
\planteamiento por el \cref{teo:duales}, la parte dual de $f(a+\varepsilon)$ es $f'(a)$; por el \cref{cor:coste-directo}, cada pasada da una derivada direccional.
\pasos
\begin{enumerate}[label=\arabic*.]
  \item $(3+\varepsilon)^{2}=9+6\varepsilon+\varepsilon^{2}=9+6\varepsilon$.
  \item La parte dual, 6, es la derivada de $x^{2}$ en $x=3$.
  \item El gradiente de $f\colon\R^{1000}\to\R$ tiene $1000$ componentes, y cada pasada en modo directo con semilla $e_k$ da una: $1000$ pasadas.
  \item Una derivada direccional $\nabla f\T v$ es exactamente lo que da una pasada con semilla $v$: una pasada.
\end{enumerate}
\resultado \textbf{(a) parte dual 6, la derivada de $x^{2}$ en 3; (b) 1000 pasadas para el gradiente y una para una derivada direccional.}
\comprobacion $\frac{d}{dx}x^{2}=2x=6$ en $x=3$.
\end{solucion}

\begin{solucion}{eje:2-cuartica}
\datos $f(x)=x^{4}-3x$, $x=1$.
\planteamiento evaluar $f(1+\varepsilon)$ con $\varepsilon^{2}=0$ y leer la parte real y la dual.
\pasos
\begin{enumerate}[label=\arabic*.]
  \item $x^{2}=(1+\varepsilon)(1+\varepsilon)=1+2\varepsilon$.
  \item $x^{4}=(1+2\varepsilon)(1+2\varepsilon)=1+4\varepsilon$.
  \item $3x=3+3\varepsilon$.
  \item $f=(1+4\varepsilon)-(3+3\varepsilon)=-2+\varepsilon$.
\end{enumerate}
\resultado \textbf{$f(1)=-2$ y $f'(1)=1$.}
\comprobacion $f'(x)=4x^{3}-3$ vale $4-3=1$ en $x=1$.
\end{solucion}

\begin{solucion}{eje:2-cociente}
\datos $g(x)=x/(1+x^{2})$, $x=2$; $1/(a+b\varepsilon)=1/a-(b/a^{2})\varepsilon$.
\planteamiento calcular el denominador como dual, invertirlo con la \cref{prop:inverso} y multiplicar por $x$.
\pasos
\begin{enumerate}[label=\arabic*.]
  \item $x^{2}=(2+\varepsilon)^{2}=4+4\varepsilon$, luego $1+x^{2}=5+4\varepsilon$.
  \item Inverso con $a=5$, $b=4$: $\tfrac15-\tfrac{4}{25}\varepsilon=0{,}2-0{,}16\varepsilon$.
  \item $(2+\varepsilon)(0{,}2-0{,}16\varepsilon)=0{,}4+\bigl(2\cdot(-0{,}16)+0{,}2\bigr)\varepsilon=0{,}4-0{,}12\varepsilon$.
\end{enumerate}
\resultado \textbf{$g(2)=2/5=0{,}4$ y $g'(2)=-3/25=-0{,}12$.}
\comprobacion $g'(x)=(1-x^{2})/(1+x^{2})^{2}$ vale $(1-4)/25=-3/25$ en $x=2$.
\end{solucion}

\begin{solucion}{eje:2-tanh}
\datos $h(x)=\tanh(x^{2})$, $x=1$; $\tanh(a+b\varepsilon)=\tanh a+(1-\tanh^{2}a)\,b\,\varepsilon$.
\planteamiento evaluar $x^{2}$ en duales y aplicar la extensión dual de $\tanh$ (\cref{def:extension}).
\pasos
\begin{enumerate}[label=\arabic*.]
  \item $x^{2}=(1+\varepsilon)^{2}=1+2\varepsilon$.
  \item $\tanh 1=0{,}76159$ y $1-\tanh^{2}1=1-0{,}58002=0{,}41997$.
  \item $\tanh(1+2\varepsilon)=0{,}76159+0{,}41997\cdot 2\,\varepsilon=0{,}7616+0{,}8399\,\varepsilon$.
\end{enumerate}
\resultado \textbf{$h(1)\approx 0{,}7616$ y $h'(1)\approx 0{,}8399$.}
\comprobacion por la regla de la cadena, $h'(x)=(1-\tanh^{2}x^{2})\cdot 2x$, que en $x=1$ vale $0{,}41997\cdot 2=0{,}8399$.
\end{solucion}

\begin{solucion}{eje:2-multidual}
\datos $x_1=2+\varepsilon_1$, $x_2=5+\varepsilon_2$, $\varepsilon_r\varepsilon_s=0$; $\cos 5=0{,}2837$.
\planteamiento cada variable lleva un valor y dos partes duales; la parte en $\varepsilon_r$ es la derivada respecto a $x_r$ (\cref{obs:multidual}).
\pasos
\begin{enumerate}[label=\arabic*.]
  \item $v_1=\ln(2+\varepsilon_1)=0{,}6931+0{,}5\,\varepsilon_1$.
  \item $v_2=(2+\varepsilon_1)(5+\varepsilon_2)=10+5\varepsilon_1+2\varepsilon_2$, porque $\varepsilon_1\varepsilon_2=0$.
  \item $v_3=\sen(5+\varepsilon_2)=-0{,}9589+0{,}2837\,\varepsilon_2$.
  \item $v_4=v_1+v_2=10{,}6931+5{,}5\,\varepsilon_1+2\,\varepsilon_2$.
  \item $v_5=v_4-v_3=11{,}6521+5{,}5\,\varepsilon_1+(2-0{,}2837)\,\varepsilon_2=11{,}6521+5{,}5\,\varepsilon_1+1{,}7163\,\varepsilon_2$.
  \item Derivada direccional en $(1,1)\T$: $\nabla f\T(1,1)\T=5{,}5+1{,}7163=7{,}2163$.
\end{enumerate}
\resultado \textbf{$\nabla f(2,5)=(5{,}5;\,1{,}7163)\T$ en una pasada; derivada direccional $7{,}2163$.}
\comprobacion coincide con el \cref{ej:tangente} y con la traza adjunta del \cref{ej:adjunta}; además, la semilla única $x=(2+\varepsilon,5+\varepsilon)$ daría directamente $7{,}2163$.
\end{solucion}

\begin{solucion}{eje:2-polinomios}
\datos polinomios reales; aritmética de la \cref{def:duales}.
\planteamiento inducción sobre el grado con la descomposición $p(x)=x\,q(x)+c$; después, la regla de la cadena para la composición.
\pasos
\begin{enumerate}[label=\arabic*.]
  \item Base (grado 0): $p=c$ constante. $p(a+b\varepsilon)=c=c+0\cdot b\,\varepsilon$ y $p'(a)=0$.
  \item Hipótesis: todo polinomio $q$ de grado a lo sumo $n$ cumple $q(a+b\varepsilon)=q(a)+q'(a)b\varepsilon$.
  \item Todo $p$ de grado $n+1$ se escribe $p(x)=x\,q(x)+c$ con $q$ de grado $n$ y $c=p(0)$.
  \item Por la hipótesis, $p(a+b\varepsilon)=(a+b\varepsilon)\bigl(q(a)+q'(a)b\varepsilon\bigr)+c$.
  \item Producto de duales: $(a+b\varepsilon)(q(a)+q'(a)b\varepsilon)=a\,q(a)+\bigl(a\,q'(a)+q(a)\bigr)b\,\varepsilon$.
  \item Por la regla del producto, $p'(x)=q(x)+x\,q'(x)$, luego $p(a+b\varepsilon)=p(a)+p'(a)\,b\,\varepsilon$.
  \item Composición: si $f$ cumple la propiedad, $f(a+b\varepsilon)=f(a)+f'(a)b\,\varepsilon$. Aplicando la de $g$ con $a'=f(a)$ y $b'=f'(a)b$: $g(f(a+b\varepsilon))=g(f(a))+g'(f(a))f'(a)\,b\,\varepsilon$.
  \item Por la regla de la cadena, $g'(f(a))f'(a)=(g\circ f)'(a)$.
\end{enumerate}
\resultado \textbf{la propiedad vale para todo polinomio y se conserva por composición}; es el \cref{teo:duales} para funciones de una variable.
\comprobacion $p(x)=x^{3}$ en $a=2$, $b=1$: el \cref{ej:dual-cubo} dio $(2+\varepsilon)^{3}=8+12\varepsilon$ y $p'(2)=12$.
\end{solucion}

% ---------------------------------------------------------------------
\section{Tema 3}

\begin{solucion}{eje:3-acumular}
\datos actualización $\bar v_j\mathrel{+}=\bar v_i\,\partial\varphi_i/\partial v_j$; ejemplo $y=x\cdot x$ con $x=3$.
\planteamiento por el \cref{teo:inverso}, el adjunto de una variable es la suma de las contribuciones de todas las operaciones que la usan.
\pasos
\begin{enumerate}[label=\arabic*.]
  \item El producto $y=a\cdot b$ con $a=b=x$ ejecuta $\bar a\mathrel{+}=b\,\bar y$ y $\bar b\mathrel{+}=a\,\bar y$, y ambas líneas actúan sobre $\bar x$.
  \item Con acumulación: $\bar x=0+3\cdot 1+3\cdot 1=6$.
  \item Con asignación: la primera línea pone $\bar x=3$ y la segunda lo sobrescribe con $\bar x=3$.
\end{enumerate}
\resultado \textbf{con acumulación se obtiene $2x=6$, correcto; con asignación, $3$, erróneo.}
\comprobacion $\frac{d}{dx}x^{2}=6$ en $x=3$. El \cref{eje:3-arbol} caracteriza cuándo la asignación basta.
\end{solucion}

\begin{solucion}{eje:3-producto}
\datos $f(x,y)=(x+y)\cdot(x\cdot y)$ en $(2,3)$.
\planteamiento nodos $s=x+y$, $p=x\cdot y$, $f=s\cdot p$; $x$ e $y$ alimentan a dos nodos y sus adjuntos acumulan dos contribuciones.
\pasos
\begin{enumerate}[label=\arabic*.]
  \item Hacia delante: $s=5$, $p=6$, $f=30$.
  \item $\bar f=1$. Nodo $f=s\cdot p$: $\bar s=p=6$, $\bar p=s=5$.
  \item Nodo $s=x+y$: $\bar x\mathrel{+}=6$, $\bar y\mathrel{+}=6$.
  \item Nodo $p=x\cdot y$: $\bar x\mathrel{+}=\bar p\,y=15$, luego $\bar x=21$; $\bar y\mathrel{+}=\bar p\,x=10$, luego $\bar y=16$.
\end{enumerate}
\resultado \textbf{$f=30$, $\partial f/\partial x=21$, $\partial f/\partial y=16$.}
\comprobacion $f=x^{2}y+xy^{2}$: $\partial_xf=2xy+y^{2}=12+9=21$ y $\partial_yf=x^{2}+2xy=4+12=16$.
\end{solucion}

\begin{solucion}{eje:3-neurona}
\datos $o=\tanh(n)$, $n=w_1x_1+w_2x_2+b$, $x=(2,0)$, $w=(-1,3)$, $b=1$.
\planteamiento como en el \cref{ej:neurona}: $\bar n=(1-o^{2})\bar o$ con $\bar o=1$, y después las reglas de la suma y del producto.
\pasos
\begin{enumerate}[label=\arabic*.]
  \item $n=(-1)\cdot 2+3\cdot 0+1=-1$.
  \item $o=\tanh(-1)=-0{,}76159$; $1-o^{2}=0{,}41997=\bar n$.
  \item Suma: $\bar b=\bar n=0{,}4200$.
  \item Productos: $\bar w_1=\bar n\,x_1=0{,}8399$; $\bar w_2=\bar n\,x_2=0$; $\bar x_1=\bar n\,w_1=-0{,}4200$; $\bar x_2=\bar n\,w_2=1{,}260$.
\end{enumerate}
\resultado \textbf{$o\approx-0{,}7616$; $\partial o/\partial w_1\approx 0{,}8399$, $\partial o/\partial w_2=0$, $\partial o/\partial b\approx 0{,}4200$, $\partial o/\partial x_1\approx-0{,}4200$, $\partial o/\partial x_2\approx 1{,}260$.}
\comprobacion un peso cuya entrada es nula no recibe gradiente, y $\partial o/\partial w_1=x_1\,\partial o/\partial b=2\cdot 0{,}41997$.
\end{solucion}

\begin{solucion}{eje:3-topologico}
\datos $f=a\cdot b+a$, $a=2$, $b=5$; nodos $a$, $b$, $m=a\cdot b$, $f=m+a$.
\planteamiento $m$ debe ir tras $a$ y $b$; $f$, tras $m$ y $a$.
\pasos
\begin{enumerate}[label=\arabic*.]
  \item Órdenes topológicos: solo pueden variar las posiciones relativas de $a$ y $b$, que van antes que $m$, y $f$ va al final: $(a,b,m,f)$ y $(b,a,m,f)$.
  \item Hacia delante: $m=10$, $f=12$.
  \item $\bar f=1$. Nodo $f=m+a$: $\bar m=1$, $\bar a\mathrel{+}=1$.
  \item Nodo $m=a\cdot b$: $\bar a\mathrel{+}=\bar m\,b=5$, luego $\bar a=6$; $\bar b\mathrel{+}=\bar m\,a=2$.
\end{enumerate}
\resultado \textbf{dos órdenes topológicos; $f=12$, $\partial f/\partial a=6$, $\partial f/\partial b=2$.}
\comprobacion $\partial f/\partial a=b+1=6$ y $\partial f/\partial b=a=2$.
\end{solucion}

\begin{solucion}{eje:3-d2l}
\datos el grafo del \cref{eje:1-d2l}: $a=x^{2}$, $b=\ln a$, $c=\sen x$, $d=b\cdot c$, $e=1/x$, $f=d+e$, en $x=2$, con $a=4$, $b=1{,}386$, $c=0{,}9093$.
\planteamiento modo directo con $\dot x=1$; modo inverso con $\bar f=1$, sumando en $\bar x$ las tres contribuciones de $a$, $c$ y $e$.
\pasos
\begin{enumerate}[label=\arabic*.]
  \item (a) Derivadas locales: $\partial a/\partial x=2x=4$; $\partial b/\partial a=1/a=0{,}25$; $\partial c/\partial x=\cos 2=-0{,}4161$; $\partial d/\partial b=c=0{,}9093$; $\partial d/\partial c=b=1{,}386$; $\partial e/\partial x=-1/x^{2}=-0{,}25$; $\partial f/\partial d=\partial f/\partial e=1$.
  \item (b) Directo: $\dot a=4$; $\dot b=\dot a/a=1$; $\dot c=-0{,}4161$; $\dot d=\dot b\,c+b\,\dot c=0{,}9093-0{,}5769=0{,}3324$; $\dot e=-0{,}25$; $\dot f=0{,}3324-0{,}25=0{,}0824$.
  \item (c) Inverso: $\bar d=\bar e=1$; $\bar b=\bar d\,c=0{,}9093$; $\bar c=\bar d\,b=1{,}386$; $\bar a=\bar b/a=0{,}2273$.
  \item Contribuciones a $\bar x$: desde $a$, $0{,}2273\cdot 4=0{,}9093$; desde $c$, $1{,}386\cdot(-0{,}4161)=-0{,}5769$; desde $e$, $1\cdot(-0{,}25)=-0{,}25$. Suma: $\bar x=0{,}0824$.
\end{enumerate}
\resultado \textbf{$f'(2)\approx 0{,}0824$ por los dos modos.}
\comprobacion $f'(x)=\frac{2}{x}\sen x+\ln(x^{2})\cos x-\frac{1}{x^{2}}$; en $x=2$: $0{,}9093-0{,}5769-0{,}25=0{,}0824$. PyTorch da $0{,}08240$.
\end{solucion}

\begin{solucion}{eje:3-sesgo}
\datos $x=(1,2)$, $W_1=\begin{psmallmatrix}1&0\\-1&0{,}5\end{psmallmatrix}$, $b_1=(0{,}5;\,0)$, $W_2=\begin{psmallmatrix}2\\1\end{psmallmatrix}$, $y=3$, $\lambda=0{,}1$.
\planteamiento el sesgo es un sumando más de $z$, luego $\bar b_1=\bar z$; no aparece en $s$ y no recibe término de regularización. Cada peso recibe la parte de $L$ más $\lambda W$, por la arista que lo une a $s$.
\pasos
\begin{enumerate}[label=\arabic*.]
  \item (a) Ecuaciones: $\bar o=o-y$; $\bar W_2=h\T\bar o+\lambda W_2$; $\bar h=\bar o\,W_2\T$; $\bar z=\bar h\odot\ind[z>0]$; $\bar W_1=x\T\bar z+\lambda W_1$; $\bar b_1=\bar z$.
  \item (b) $xW_1=(1-2,\;0+1)=(-1,1)$; $z=(-0{,}5;\,1)$; $h=(0,1)$.
  \item $o=hW_2=1$; $L=\tfrac12(1-3)^{2}=2$.
  \item $\norm{W_1}_F^{2}=1+0+1+0{,}25=2{,}25$; $\norm{W_2}_F^{2}=5$; $s=0{,}05\cdot 7{,}25=0{,}3625$; $J=2{,}3625$.
  \item $\bar o=-2$; $\bar W_2=(0;\,-2)\T+(0{,}2;\,0{,}1)\T=(0{,}2;\,-1{,}9)\T$.
  \item $\bar h=-2\cdot(2,1)=(-4,-2)$; como $z_1<0$, $\bar z=(0,-2)$.
  \item $x\T\bar z=\begin{psmallmatrix}0&-2\\0&-4\end{psmallmatrix}$ y $\bar W_1=\begin{psmallmatrix}0&-2\\0&-4\end{psmallmatrix}+\begin{psmallmatrix}0{,}1&0\\-0{,}1&0{,}05\end{psmallmatrix}=\begin{psmallmatrix}0{,}1&-2\\-0{,}1&-3{,}95\end{psmallmatrix}$.
  \item $\bar b_1=\bar z=(0,-2)$.
\end{enumerate}
\resultado \textbf{$J=2{,}3625$; $\nabla_{\!W_2}J=(0{,}2;\,-1{,}9)\T$; $\nabla_{\!W_1}J=\begin{psmallmatrix}0{,}1&-2\\-0{,}1&-3{,}95\end{psmallmatrix}$; $\nabla_{\!b_1}J=(0,-2)$.}
\comprobacion la primera columna de $\nabla_{\!W_1}J$ es solo $\lambda W_1$, porque la unidad~1 está inactiva. PyTorch da los mismos valores.
\end{solucion}

\begin{solucion}{eje:3-arbol}
\datos el algoritmo de la \cref{fig:algoritmo} con asignación en lugar de acumulación.
\planteamiento si cada variable es argumento de una sola operación, cada acumulador recibe una única contribución, y sumarla a~0 equivale a asignarla.
\pasos
\begin{enumerate}[label=\arabic*.]
  \item Sea $j<\ell$. Si $v_j$ es argumento de una única operación $i$, y figura una sola vez en su lista de argumentos, la versión con acumulación ejecuta sobre $\bar v_j$ una sola actualización, $\bar v_j\gets 0+\bar v_i\,\partial\varphi_i/\partial v_j$, que coincide con la asignación.
  \item La salida $v_\ell$ no es argumento de ninguna operación y conserva la semilla $\bar v_\ell=1$ en ambas versiones.
  \item Por tanto, las dos versiones ejecutan las mismas operaciones con los mismos valores, y por el \cref{teo:inverso} la versión con asignación es correcta.
  \item Contraejemplo con un nodo de dos sucesores: $u=2x$, $z=u^{2}+u$ con $x=3$, es decir, $y=u^{2}$ y $z=y+u$. Con asignación, el nodo $z$ pone $\bar y=1$ y $\bar u=1$; el nodo $y$ \emph{sustituye} $\bar u$ por $2u\cdot\bar y=12$, y el nodo $u$ da $\bar x=2\cdot 12=24$.
  \item El valor correcto es $dz/dx=8x+2=26$: la asignación ha perdido la contribución $\bar u=1$ del camino directo.
\end{enumerate}
\resultado \textbf{la asignación es correcta en un árbol; con un nodo usado dos veces puede fallar}, como en el ejemplo, que da $24$ en lugar de $26$.
\comprobacion con acumulación, $\bar u=1+12=13$ y $\bar x=26$ (\cref{ej:orden}).
\end{solucion}

% ---------------------------------------------------------------------
\section{Tema 4}

\begin{solucion}{eje:4-modo}
\datos $f\colon\R^{1000}\to\R$ ($n=1000$, $m=1$); $g\colon\R\to\R^{1000}$ ($n=1$, $m=1000$).
\planteamiento el modo directo da una columna del jacobiano por pasada y el inverso una fila (\cref{cor:eleccion}).
\pasos
\begin{enumerate}[label=\arabic*.]
  \item $J_f$ es $1\times 1000$, una fila: el modo inverso necesita 1 pasada; el directo, 1000.
  \item $J_g$ es $1000\times 1$, una columna: el modo directo necesita 1 pasada; el inverso, 1000.
\end{enumerate}
\resultado \textbf{$f$: modo inverso, 1 pasada; $g$: modo directo, 1 pasada.} El modo contrario costaría 1000 pasadas en ambos casos.
\comprobacion el número de pasadas es $\min(n,m)=1$ en los dos casos, eligiendo el modo adecuado.
\end{solucion}

\begin{solucion}{eje:4-d2l}
\datos los dos modos de la AD aplicados a $f\colon\R^{n}\to\R^{m}$.
\planteamiento comparar los modos según los tres criterios del enunciado, usando las \cref{cor:coste-directo,prop:vjp,prop:memoria} y la \cref{prop:asociatividad}.
\pasos
\begin{enumerate}[label=\arabic*.]
  \item \emph{Datos intermedios.} El modo directo descarta cada tangente en cuanto deja de usarse y necesita memoria comparable a la de evaluar $f$. El inverso debe guardar los valores de toda la pasada hacia delante (\cref{prop:memoria}).
  \item \emph{Paralelismo.} En el modo directo, las $n$ pasadas con semillas distintas son independientes y pueden ejecutarse en paralelo, o en una sola pasada vectorial (\cref{obs:multidual}). En el inverso, la pasada hacia atrás no puede empezar hasta terminar la pasada hacia delante.
  \item \emph{Tamaños.} El directo propaga matrices $m'\times n$ o vectores de tamaño $n$ por dirección; el inverso propaga vectores fila del tamaño de cada variable intermedia. Con una salida escalar y muchas entradas, el inverso multiplica siempre vectores por matrices (\cref{prop:asociatividad}).
\end{enumerate}
\resultado \textbf{directo cuando hay pocas entradas ($n\ll m$) o la memoria es crítica; inverso cuando hay pocas salidas ($m\ll n$), como en el entrenamiento, con una pérdida escalar.}
\comprobacion un LLM tiene $n\approx 10^{11}$ y $m=1$: el inverso cuesta una pasada, el directo $10^{11}$.
\end{solucion}

\begin{solucion}{eje:4-cadena}
\datos $c\in\R^{1\times 100}$; $A,B\in\R^{100\times 100}$.
\planteamiento un producto $(p\times q)(q\times r)$ cuesta $pqr$ multiplicaciones (\cref{prop:asociatividad} con $L=3$).
\pasos
\begin{enumerate}[label=\arabic*.]
  \item $c(BA)$: $BA$ cuesta $100^{3}=1\,000\,000$; $c\,(BA)$ cuesta $100^{2}=10\,000$. Total: $1\,010\,000$.
  \item $(cB)A$: $cB$ cuesta $10\,000$ y da un vector $1\times 100$; $(cB)A$ cuesta otros $10\,000$. Total: $20\,000$.
  \item Cociente: $1\,010\,000/20\,000=50{,}5$.
  \item En general: $(m^{3}+m^{2})/(2m^{2})=(m+1)/2$.
\end{enumerate}
\resultado \textbf{$1\,010\,000$ frente a $20\,000$ multiplicaciones; cociente $50{,}5=(m+1)/2$ con $m=100$.}
\comprobacion con $m=1000$ el cociente es $1001/2=500{,}5$: crece linealmente con $m$.
\end{solucion}

\begin{solucion}{eje:4-lineal}
\datos $X=(2,-1)$, $W=\begin{psmallmatrix}1&2&0\\0&1&-1\end{psmallmatrix}$, $\bar Y=(1,0,2)$.
\planteamiento \cref{prop:vjp-lineal}: $\bar X=\bar YW\T$ y $\bar W=X\T\bar Y$.
\pasos
\begin{enumerate}[label=\arabic*.]
  \item $Y=XW=(2\cdot 1+(-1)\cdot 0,\;2\cdot 2+(-1)\cdot 1,\;2\cdot 0+(-1)\cdot(-1))=(2,3,1)$.
  \item $W\T=\begin{psmallmatrix}1&0\\2&1\\0&-1\end{psmallmatrix}$; $\bar X=\bar YW\T=(1\cdot 1+0\cdot 2+2\cdot 0,\;1\cdot 0+0\cdot 1+2\cdot(-1))=(1,-2)$.
  \item $\bar W=X\T\bar Y=\begin{psmallmatrix}2\\-1\end{psmallmatrix}(1,0,2)=\begin{psmallmatrix}2&0&4\\-1&0&-2\end{psmallmatrix}$.
  \item Formas: $\bar X$ es $1\times 2$, como $X$; $\bar W$ es $2\times 3$, como $W$.
  \item Comprobación directa: $L=\bar YY\T=Y_1+2Y_3$ y $Y_3=X_1W_{13}+X_2W_{23}$, luego $\partial L/\partial W_{13}=2X_1=4$.
\end{enumerate}
\resultado \textbf{$Y=(2,3,1)$; $\bar X=(1,-2)$; $\bar W=\begin{psmallmatrix}2&0&4\\-1&0&-2\end{psmallmatrix}$.}
\comprobacion la columna central de $\bar W$ es nula porque $\bar Y_2=0$, y la entrada $(1,3)$ coincide con la derivada directa.
\end{solucion}

\begin{solucion}{eje:4-cs336}
\datos $N=10^{9}$; $512\times 1024$ tokens por paso; $10^{5}$ pasos; pico $495\cdot 10^{12}$ FLOP/s; utilización $0{,}5$.
\planteamiento FLOPs por paso $=(2N+4N)\cdot$ tokens por paso (\cref{cor:6nd}); tiempo $=$ FLOPs totales$/$FLOP/s efectivos.
\pasos
\begin{enumerate}[label=\arabic*.]
  \item Tokens por paso: $512\cdot 1024=524\,288$.
  \item FLOPs por paso: $6\cdot 10^{9}\cdot 524\,288=3{,}146\cdot 10^{15}$.
  \item Total: $3{,}146\cdot 10^{15}\cdot 10^{5}=3{,}146\cdot 10^{20}$ FLOPs.
  \item Rendimiento efectivo: $0{,}5\cdot 495\cdot 10^{12}=2{,}475\cdot 10^{14}$ FLOP/s.
  \item Tiempo: $3{,}146\cdot 10^{20}/2{,}475\cdot 10^{14}=1{,}271\cdot 10^{6}$~s, es decir, $1{,}271\cdot 10^{6}/3600=353{,}1$~h y $353{,}1/24=14{,}71$ días.
\end{enumerate}
\resultado \textbf{$\approx 3{,}15\cdot 10^{20}$ FLOPs; $\approx 353$~h, unos $14{,}7$ días en una GPU.}
\comprobacion $D=524\,288\cdot 10^{5}\approx 5{,}24\cdot 10^{10}$ tokens y $6ND=6\cdot 10^{9}\cdot 5{,}24\cdot 10^{10}\approx 3{,}15\cdot 10^{20}$.
\end{solucion}

\begin{solucion}{eje:4-softmax}
\datos $p_i=e^{z_i}/S$, $S=\sum_ke^{z_k}$; $z=(2,1,0)$, $y=0$.
\planteamiento derivar $\ln p_i$ en lugar de $p_i$, usando $\partial p_i/\partial z_j=p_i\,\partial\ln p_i/\partial z_j$.
\pasos
\begin{enumerate}[label=\arabic*.]
  \item (a) $\ln p_i=z_i-\ln S$ da $\partial\ln p_i/\partial z_j=\ind[i=j]-e^{z_j}/S=\ind[i=j]-p_j$.
  \item Luego $\partial p_i/\partial z_j=p_i(\ind[i=j]-p_j)$ y $\bar z_j=\sum_i\bar p_ip_i(\ind[i=j]-p_j)=p_j\bar p_j-p_j\langle\bar p,p\rangle$, que es la \cref{prop:vjp-softmax}.
  \item (b) Para $L=-\ln p_y$, $\bar p=-e_y/p_y$ y $\langle\bar p,p\rangle=-1$. Entonces $\bar z_j=p_j(\bar p_j+1)$: vale $p_y-1$ si $j=y$ y $p_j$ si $j\ne y$, es decir, $\bar z=p-e_y$.
  \item Suma de las componentes: $\sum_j(p_j-\ind[j=y])=1-1=0$.
  \item (c) $e^{z}=(7{,}389;\,2{,}718;\,1)$, con $S=11{,}107$; $p=(0{,}6652;\,0{,}2447;\,0{,}0900)$.
  \item $\nabla_{\!z}L=(0{,}6652-1;\,0{,}2447;\,0{,}0900)=(-0{,}3348;\,0{,}2447;\,0{,}0900)$.
  \item $L=-\ln 0{,}6652=0{,}4076$.
\end{enumerate}
\resultado \textbf{$\nabla_{\!z}L=\softmax(z)-e_y$; con $z=(2,1,0)$ e $y=0$, $\nabla_{\!z}L\approx(-0{,}3348;\,0{,}2447;\,0{,}0900)$ y $L\approx 0{,}4076$.}
\comprobacion $-0{,}3348+0{,}2447+0{,}0900=-0{,}0001\approx 0$, por redondeo; también $L=-2+\ln 11{,}107=0{,}4076$. Coincide con el \cref{ej:ce-micro} y con \texttt{F.cross\_entropy}.
\end{solucion}

% ---------------------------------------------------------------------
\section{Tema 5}

\begin{solucion}{eje:5-d2l}
\datos $f\colon\R^{n}\to\R$; gradiente con un coste $O(\ops(f))$ (\cref{teo:barato}).
\planteamiento el hessiano tiene $n^{2}$ entradas, frente a las $n$ del gradiente, y cada columna se obtiene con un producto hessiana-vector (\cref{prop:hvp}).
\pasos
\begin{enumerate}[label=\arabic*.]
  \item La segunda derivada se obtiene aplicando la AD al programa que calcula el gradiente, que ya es una o dos veces más largo que el de $f$ y contiene las derivadas locales.
  \item Un producto $H_fv$ cuesta, por la \cref{prop:hvp}, una constante por $\ops(f)$: del mismo orden que el gradiente.
  \item El hessiano completo necesita $n$ productos hessiana-vector, uno por columna: coste del orden de $n\cdot\ops(f)$, y $n^{2}$ números que guardar.
  \item Además, el programa del gradiente debe conservar su propio grafo para poder derivarlo de nuevo, lo que aumenta la memoria (en PyTorch, \texttt{create\_graph=True}).
\end{enumerate}
\resultado \textbf{un producto hessiana-vector cuesta lo mismo que un gradiente, salvo una constante; el hessiano completo cuesta $n$ veces más y ocupa $n^{2}$ números.}
\comprobacion con $n=10^{9}$, el hessiano tendría $10^{18}$ entradas, inabordable; los métodos que usan curvatura emplean productos hessiana-vector.
\end{solucion}

\begin{solucion}{eje:5-hvp}
\datos $f=x_1^{2}x_2+x_2^{3}$, punto $(1,2)$, $v=(0,1)\T$.
\planteamiento calcular $\nabla f=(2x_1x_2,\;x_1^{2}+3x_2^{2})$ con $x_1=1$ y $x_2=2+\varepsilon$; la parte dual es $Hv$ (\cref{prop:hvp}).
\pasos
\begin{enumerate}[label=\arabic*.]
  \item Primera componente: $2\cdot 1\cdot(2+\varepsilon)=4+2\varepsilon$.
  \item $(2+\varepsilon)^{2}=4+4\varepsilon$, luego $3x_2^{2}=12+12\varepsilon$.
  \item Segunda componente: $1+12+12\varepsilon=13+12\varepsilon$.
  \item Lectura: $\nabla f(1,2)=(4,13)$ y $Hv=(2,12)$.
  \item Hessiano: $H=\begin{psmallmatrix}2x_2&2x_1\\2x_1&6x_2\end{psmallmatrix}=\begin{psmallmatrix}4&2\\2&12\end{psmallmatrix}$, cuya segunda columna es $(2,12)$.
\end{enumerate}
\resultado \textbf{$\nabla f(1,2)=(4,13)$ y $H_f(1,2)\,v=(2,12)$.}
\comprobacion PyTorch, derivando el gradiente con \texttt{create\_graph=True}, da los mismos vectores.
\end{solucion}

\begin{solucion}{eje:5-adjunto}
\datos $x_{t+1}=ax_t+b$, $t=0,1,2$; $x_0=1$, $a=2$, $b=1$; $L=x_3$.
\planteamiento por el \cref{teo:lagrange}, $\lambda_3=\partial L/\partial x_3=1$, $\lambda_t=\lambda_{t+1}\,\partial x_{t+1}/\partial x_t=a\lambda_{t+1}$, y las derivadas respecto a los parámetros suman las contribuciones de cada transición.
\pasos
\begin{enumerate}[label=\arabic*.]
  \item Trayectoria: $x_1=3$, $x_2=7$, $x_3=15$.
  \item Multiplicadores: $\lambda_3=1$, $\lambda_2=2$, $\lambda_1=4$.
  \item $dL/da=\sum_{t=0}^{2}\lambda_{t+1}\,x_t=\lambda_1x_0+\lambda_2x_1+\lambda_3x_2=4+6+7=17$.
  \item $dL/db=\sum_{t=0}^{2}\lambda_{t+1}\cdot 1=4+2+1=7$.
  \item Expresión cerrada: $x_3=a^{3}x_0+b(a^{2}+a+1)$.
\end{enumerate}
\resultado \textbf{$x=(1,3,7,15)$; $\lambda=(\lambda_1,\lambda_2,\lambda_3)=(4,2,1)$; $dL/da=17$ y $dL/db=7$.}
\comprobacion $dx_3/da=3a^{2}x_0+b(2a+1)=12+5=17$ y $dx_3/db=a^{2}+a+1=7$.
\end{solucion}

\begin{solucion}{eje:5-checkpoint}
\datos memoria $M(k)=L/k+k$ en unidades de activación; $L=16$ y $L=100$.
\planteamiento minimizar $M$ con derivadas y contar las pasadas como en la \cref{prop:flops}.
\pasos
\begin{enumerate}[label=\arabic*.]
  \item Se guardan $L/k$ puntos de control; al retropropagar un tramo se recalculan y guardan sus $k$ activaciones. La memoria máxima es $M(k)=L/k+k$.
  \item $M'(k)=1-L/k^{2}=0\iff k=\sqrt L$ (con $k>0$); $M''(k)=2L/k^{3}>0$: mínimo, de valor $2\sqrt L$.
  \item $L=16$: $k=4$, $M=4+4=8$, frente a 16 sin puntos de control.
  \item $L=100$: $k=10$, $M=10+10=20$, frente a 100.
  \item Cómputo: pasada hacia delante ($2N$ FLOPs por token), recálculo ($2N$) y pasada hacia atrás ($4N$): $8N$ por token, un $8/6-1=33\,\%$ más.
\end{enumerate}
\resultado \textbf{$k=\sqrt L$ y memoria $2\sqrt L$: 8 para $L=16$ y 20 para $L=100$; el paso cuesta $8N$ FLOPs por token.}
\comprobacion con $L=16$, $k=2$ da $8+2=10$ y $k=8$ da $2+8=10$, ambos mayores que 8.
\end{solucion}

\begin{solucion}{eje:5-cs336}
\datos $N_b=16$; coste~1 por bloque de activaciones y $c=1/4$ por punto de control; $M(k)=(N_b/k)c+k$.
\planteamiento equilibrar la memoria de los puntos de control, decreciente en $k$, con la del tramo recalculado, creciente, como en el \cref{eje:5-checkpoint} pero con $c\ne 1$.
\pasos
\begin{enumerate}[label=\arabic*.]
  \item (a) $k=1$: $16\cdot 0{,}25+1=5$. $k=2$: $8\cdot 0{,}25+2=4$. $k=4$: $4\cdot 0{,}25+4=5$. $k=8$: $2\cdot 0{,}25+8=8{,}5$. $k=16$: $0{,}25+16=16{,}25$.
  \item Óptimo continuo: $M'(k)=-N_bc/k^{2}+1=0\iff k=\sqrt{N_bc}=\sqrt 4=2$, con $M(2)=4$, frente a 16 sin recálculo.
  \item (b) Estrategia recursiva: se parte la pila en dos mitades con un punto de control al inicio de cada una; para retropropagar una mitad, se vuelve a partir en dos, y así hasta bloques sueltos. Hay $\log_2N_b$ niveles.
  \item Memoria: en cada nivel se mantiene un número acotado de puntos de control, más las activaciones de un solo bloque: $O(\log N_b)$; con $N_b=16$, unos $\log_216=4$ niveles.
  \item Cómputo: cada bloque se recalcula una vez por nivel, $O(\log N_b)$ veces, y el total es de $O(N_b\log N_b)$ evaluaciones de bloque, frente a $O(N_b)$ sin recálculo.
\end{enumerate}
\resultado \textbf{(a) $M=5;\,4;\,5;\,8{,}5;\,16{,}25$, con óptimo $k=\sqrt{N_bc}=2$ y $M=4$. (b) Recálculo recursivo por mitades: memoria $O(\log N_b)$, cómputo $O(N_b\log N_b)$.}
\comprobacion con $c=1$ la fórmula se reduce a la del \cref{eje:5-checkpoint}, con óptimo $k=\sqrt{N_b}=4$.
\end{solucion}

% ---------------------------------------------------------------------
\section{Tema 6}

\begin{solucion}{eje:6-zero}
\datos bucle de PyTorch sin \texttt{opt.zero\_grad()}; evaluación con \texttt{torch.no\_grad()}.
\planteamiento \texttt{backward()} suma en \texttt{.grad} (acumulación del \cref{tema:3}); \texttt{no\_grad} impide registrar el grafo.
\pasos
\begin{enumerate}[label=\arabic*.]
  \item Sin \texttt{zero\_grad()}, en el paso $t$ el gradiente guardado es $g_1+\dots+g_t$, no $g_t$.
  \item Si los $g_i$ son parecidos, el paso efectivo crece como $t\,\eta\,g$: equivale a multiplicar la tasa de aprendizaje por $t$, y el entrenamiento se vuelve errático o diverge.
  \item \texttt{torch.no\_grad()} no crea nodos ni guarda tensores. Al evaluar no hace falta derivar, así que ahorra memoria y tiempo, y la evaluación no deja rastro en \texttt{.grad} ni en el grafo del paso siguiente.
\end{enumerate}
\resultado \textbf{sin \texttt{zero\_grad()} los gradientes se acumulan entre pasos; \texttt{no\_grad()} evalúa sin registrar el grafo.}
\comprobacion en el \cref{eje:6-dos-veces}, dos llamadas a \texttt{backward()} sobre $y=x^{2}$ en $x=3$ dejan \texttt{x.grad}${}=6+6=12$.
\end{solucion}

\begin{solucion}{eje:6-dos-veces}
\datos $x=3$, $y=x\cdot x$, dos llamadas a \texttt{backward()}.
\planteamiento PyTorch libera los tensores guardados tras la pasada hacia atrás, salvo con \texttt{retain\_graph=True}; \texttt{.grad} se acumula.
\pasos
\begin{enumerate}[label=\arabic*.]
  \item (a) Primera llamada: \texttt{x.grad}${}=0+2\cdot 3=6$, y se liberan los valores guardados del nodo producto.
  \item Segunda llamada: el VJP del producto necesita esos valores, que ya no existen, y PyTorch lanza un \texttt{RuntimeError} que indica que se intenta retropropagar por el grafo una segunda vez.
  \item (b) Con \texttt{retain\_graph=True} el grafo se conserva: tras la primera llamada \texttt{x.grad}${}=6$ y tras la segunda \texttt{x.grad}${}=6+6=12$.
\end{enumerate}
\resultado \textbf{(a) error en la segunda llamada; (b) \texttt{x.grad}${}=12$}, el doble del gradiente, por la acumulación.
\comprobacion es el mismo efecto que olvidar \texttt{zero\_grad()} (\cref{eje:6-zero}). PyTorch~2.14 reproduce ambos comportamientos.
\end{solucion}

\begin{solucion}{eje:6-cmu}
\datos operaciones $a-b$, $a\cdot b$, $a/b$, $\ln a$, $\exp a$, $AB$; $f=\ln(a/b)-\exp(a\cdot b)$ en $(2,1)$.
\planteamiento cada función hacia atrás aplica la regla local (VJP) y devuelve un adjunto por entrada; la cinta se recorre al revés (\cref{obs:cinta}).
\pasos
\begin{enumerate}[label=\arabic*.]
  \item Reglas: $a-b$: $(g,\,-g)$. $a\cdot b$: $(g\,b,\;g\,a)$. $a/b$: $(g/b,\;-g\,a/b^{2})$. $\ln a$: $g/a$. $\exp a$: $g\,e^{a}$. $AB$: $(G\,B\T,\;A\T G)$ (\cref{prop:vjp-lineal}).
  \item Cinta hacia delante: $v_1=a/b=2$; $v_2=\ln v_1=0{,}6931$; $v_3=a\cdot b=2$; $v_4=\exp v_3=7{,}389$; $f=v_2-v_4=-6{,}696$.
  \item Hacia atrás: resta, $\bar v_2=1$, $\bar v_4=-1$; exponencial, $\bar v_3=-1\cdot e^{2}=-7{,}389$; logaritmo, $\bar v_1=1/2=0{,}5$.
  \item Producto $v_3=a\cdot b$: $\bar a\mathrel{+}=-7{,}389\cdot 1$ y $\bar b\mathrel{+}=-7{,}389\cdot 2=-14{,}78$.
  \item Cociente $v_1=a/b$: $\bar a\mathrel{+}=0{,}5/1$, luego $\bar a=-6{,}889$; $\bar b\mathrel{+}=-0{,}5\cdot 2/1^{2}=-1$, luego $\bar b=-15{,}78$.
\end{enumerate}
\resultado \textbf{$f(2,1)\approx-6{,}696$, $\partial f/\partial a\approx-6{,}889$ y $\partial f/\partial b\approx-15{,}78$.}
\comprobacion $f=\ln a-\ln b-e^{ab}$: $\partial_af=1/a-b\,e^{ab}=0{,}5-7{,}389=-6{,}889$ y $\partial_bf=-1/b-a\,e^{ab}=-1-14{,}78=-15{,}78$.
\end{solucion}

\begin{solucion}{eje:6-parametros}
\datos $V=28$, $d=8$, $B=3$, $H=64$.
\planteamiento embeddings: $V\cdot d$; capa lineal de $n$ entradas y $k$ salidas: $nk$ pesos y $k$ sesgos.
\pasos
\begin{enumerate}[label=\arabic*.]
  \item Embeddings: $28\cdot 8=224$.
  \item Capa oculta: entrada $B\cdot d=24$; $24\cdot 64+64=1536+64=1600$.
  \item Capa de salida: $64\cdot 28+28=1792+28=1820$.
  \item Total: $224+1600+1820=3644$.
\end{enumerate}
\resultado \textbf{$224+1600+1820=3644$ parámetros.}
\comprobacion con $d=10$ se obtiene $280+1984+1820=4084$, la cifra del \cref{ej:mlp-caracteres}.
\end{solucion}

\begin{solucion}{eje:6-paso}
\datos $o=\tanh(wx+b)$, $x=1$, $w=0{,}5$, $b=0$, $y=1$, $L=\tfrac12(o-y)^{2}$, $\eta=0{,}5$.
\planteamiento modo inverso por los nodos $n=wx+b$, $o=\tanh n$ y $L$; después, $w\gets w-\eta\,\bar w$ y $b\gets b-\eta\,\bar b$.
\pasos
\begin{enumerate}[label=\arabic*.]
  \item $n=0{,}5$; $o=0{,}46212$; $L=\tfrac12(0{,}46212-1)^{2}=0{,}14466$.
  \item $\bar o=o-y=-0{,}53788$; $\bar n=\bar o\,(1-o^{2})=-0{,}53788\cdot 0{,}78645=-0{,}42302$.
  \item $\bar w=\bar n\,x=-0{,}4230$; $\bar b=\bar n=-0{,}4230$.
  \item Paso: $w=0{,}5+0{,}5\cdot 0{,}42302=0{,}71151$; $b=0+0{,}21151=0{,}21151$.
  \item Nueva pérdida: $n=0{,}92302$, $o=\tanh 0{,}92302=0{,}72732$ y $L=\tfrac12(0{,}72732-1)^{2}=0{,}03718$.
\end{enumerate}
\resultado \textbf{$L=0{,}1447$, $\partial L/\partial w=\partial L/\partial b\approx-0{,}4230$; tras el paso, $w\approx 0{,}7115$, $b\approx 0{,}2115$ y $L\approx 0{,}0372$.}
\comprobacion la pérdida disminuye de $0{,}1447$ a $0{,}0372$. El mismo paso en PyTorch, con \texttt{torch.no\_grad()} en la actualización, da los mismos números.
\end{solucion}

\begin{solucion}{eje:6-microlotes}
\datos $\ell_s(w)=\tfrac12(wx_s-y_s)^{2}$; $x=(1,2,3,4)$, $y=(2,3,7,8)$, $w=1$; $k=2$ microlotes, $S_1=\{1,2\}$ y $S_2=\{3,4\}$.
\planteamiento linealidad de la derivada (\cref{prop:microlotes}); $\ell_s'(w)=(wx_s-y_s)x_s$.
\pasos
\begin{enumerate}[label=\arabic*.]
  \item (a) Tras las $k$ llamadas, el acumulador vale $\sum_r\nabla(\mathcal L_r/k)=\nabla\bigl(\frac1k\sum_r\mathcal L_r\bigr)$, y con $|S_r|=|S|/k$ se tiene $\frac1k\sum_r\mathcal L_r=\mathcal L$.
  \item Sin dividir por $k$, el acumulador vale $\sum_r\nabla\mathcal L_r=k\,\nabla\mathcal L$: el gradiente queda multiplicado por $k$, como si la tasa de aprendizaje fuera $k$ veces mayor.
  \item (b) Residuos $wx_s-y_s=(-1,-1,-4,-4)$; derivadas $\ell_s'=(-1,-2,-12,-16)$.
  \item Lote completo: $\nabla\mathcal L=(-1-2-12-16)/4=-31/4=-7{,}75$.
  \item Microlote 1: $\nabla\mathcal L_1=(-1-2)/2=-1{,}5$; acumulado tras dividir por 2: $-0{,}75$.
  \item Microlote 2: $\nabla\mathcal L_2=(-12-16)/2=-14$; se suma $-7$: acumulado $-7{,}75$.
\end{enumerate}
\resultado \textbf{el gradiente acumulado es $-7{,}75$, igual al del lote completo}; sin dividir por $k$ sería $-15{,}5$.
\comprobacion PyTorch, con dos llamadas a \texttt{backward()} sin \texttt{zero\_grad()} intermedio, deja \texttt{w.grad}${}=-0{,}75$ y después $-7{,}75$.
\end{solucion}
